\documentclass[final,t]{beamer}
\mode<presentation>
\usetheme{Berlin}
\usepackage[english]{babel}
\usepackage[latin1]{inputenc}
\usepackage{amsmath,amsthm,amssymb,latexsym}
\usepackage{times}\usefonttheme{professionalfonts}
\boldmath
\usepackage[orientation=portrait,size=a0,scale=1.4,debug]{beamerposter}
\usepackage{graphicx}
\usepackage{tikz}
\usetikzlibrary{shapes,arrows,positioning}
\usepackage{multicol}
\usepackage{enumitem}

\title{\huge Transparent AI: The Case for Interpretability and Explainability}
\author{Dhanesh Ramachandram, Himanshu Joshi, Judy Zhu, Dhari Gandhi, Lucas Hartman, and Ananya Raval}
\institute{Vector Institute for Artificial Intelligence, Toronto}
\date{August 1, 2025}

\begin{document}
\begin{frame}[t]
\begin{columns}[t]

% Left Column
\begin{column}{.48\linewidth}

%% ABSTRACT
\begin{block}{\Large Abstract}
As AI systems increasingly inform high-stakes decisions across sectors, \textbf{transparency has become foundational} to responsible AI implementation. This paper presents:
\begin{itemize}[leftmargin=*]
\item Key insights from practical interpretability applications across diverse domains
\item Actionable strategies for organizations at varying stages of AI maturity
\item Integration of interpretability as a \textbf{core design principle} rather than retrospective add-on
\end{itemize}
\end{block}

\vspace{0.5cm}

%% MOTIVATION
\begin{block}{\Large Motivation: Why Interpretability Matters}
\textbf{Consider this scenario:}
\begin{itemize}[leftmargin=*]
\item AI diagnostic tool: 2\% error rate, no explanation
\item Human physician: 15\% error rate, fully explainable
\item \textcolor{red}{\textbf{Which would you trust?}}
\end{itemize}

\vspace{0.3cm}
\textbf{Critical Applications:}
\begin{itemize}[leftmargin=*]
\item \textbf{Healthcare:} Clinicians need to understand AI predictions to integrate with clinical judgment
\item \textbf{Finance:} Transparent credit decisions required for regulatory compliance
\item \textbf{Public Sector:} Citizens expect clear explanations for decisions affecting their lives
\end{itemize}
\end{block}

\vspace{0.5cm}

%% KEY CONCEPTS
\begin{block}{\Large Fundamentals: Interpretability vs. Explainability}

\begin{center}
\begin{tikzpicture}[scale=1.2, every node/.style={font=\Large}]
% Main boxes
\node[draw, rectangle, rounded corners, fill=blue!20, minimum width=8cm, minimum height=2.5cm, align=center] (interp) at (0,4) {\textbf{Interpretability}\\[0.3cm]``How does the model\\function internally?''};

\node[draw, rectangle, rounded corners, fill=green!20, minimum width=8cm, minimum height=2.5cm, align=center] (explain) at (0,0) {\textbf{Explainability}\\[0.3cm]``Why did the model make\\this particular decision?''};

% Arrow
\draw[->, very thick] (interp) -- (explain);

% Side annotations
\node[align=left, font=\large] at (-6,4) {Glass-box models:\\Decision trees\\Linear regression};

\node[align=left, font=\large] at (-6,0) {Post-hoc techniques:\\SHAP, LIME\\for black-box models};

\end{tikzpicture}
\end{center}

\vspace{0.3cm}
\textbf{The Trade-off Challenge:}
\begin{center}
\begin{tikzpicture}[scale=1.1]
% Axes
\draw[->, very thick] (0,0) -- (10,0) node[right, font=\Large] {Model Complexity \& Performance};
\draw[->, very thick] (0,0) -- (0,6) node[above, font=\Large] {Interpretability};

% Curve
\draw[very thick, blue] (1,5) .. controls (3,4) and (5,2.5) .. (9,1);

% Points
\filldraw[red] (1,5) circle (4pt) node[above right, font=\large] {Decision Trees};
\filldraw[red] (3,4) circle (4pt) node[above, font=\large] {Linear Models};
\filldraw[red] (5,2.5) circle (4pt) node[above, font=\large] {GAMs};
\filldraw[red] (7,1.5) circle (4pt) node[above, font=\large] {Ensemble};
\filldraw[red] (9,1) circle (4pt) node[below left, font=\large] {Deep Learning};

% Ideal region
\filldraw[green, opacity=0.3] (7,4) circle (1.5cm);
\node[font=\Large\bfseries, color=green!50!black] at (7,4) {Ideal\\Solution};

\end{tikzpicture}
\end{center}
\end{block}

\vspace{0.5cm}

%% POLICY LANDSCAPE
\begin{block}{\Large Policy \& Regulatory Landscape}
\textbf{Global AI Regulations Converge on Interpretability:}

\vspace{0.3cm}
\begin{tabular}{p{4cm}p{16cm}}
\textbf{EU AI Act} & Mandates transparency for high-risk AI systems and grants individuals right to meaningful explanations \\[0.3cm]
\textbf{Canada AIDA} & Emphasizes risk-based governance with interpretability assessments during development \\[0.3cm]
\textbf{US AI Bill of Rights} & Establishes interpretability as fundamental civil right requiring notice and explanation \\[0.3cm]
\textbf{FDA Guidance} & Requires clear documentation of AI/ML medical device decision-making \\
\end{tabular}

\vspace{0.3cm}
\colorbox{yellow!30}{\parbox{0.95\linewidth}{\large\textbf{Key Challenge:} Few regulations define what constitutes ``sufficient'' explanation or specify evaluation metrics}}
\end{block}

\end{column}

% Right Column
\begin{column}{.48\linewidth}

%% STAKEHOLDER PERSPECTIVES
\begin{block}{\Large Stakeholder-Specific Needs}

\begin{center}
\begin{tikzpicture}[scale=1.2, every node/.style={font=\large}]
% Central circle
\node[draw, circle, fill=orange!30, minimum size=3cm, align=center, font=\Large\bfseries] (center) at (0,0) {Interpretable\\AI};

% Stakeholder nodes
\node[draw, rectangle, rounded corners, fill=blue!20, minimum width=3.5cm, minimum height=1.5cm, align=center] (ds) at (-6,4) {\textbf{Data Scientists}\\Technical rigor\\Performance balance};

\node[draw, rectangle, rounded corners, fill=green!20, minimum width=3.5cm, minimum height=1.5cm, align=center] (bl) at (6,4) {\textbf{Business Leaders}\\Risk assessment\\Resource allocation};

\node[draw, rectangle, rounded corners, fill=red!20, minimum width=3.5cm, minimum height=1.5cm, align=center] (reg) at (-6,-4) {\textbf{Regulators}\\Compliance\\Standards};

\node[draw, rectangle, rounded corners, fill=purple!20, minimum width=3.5cm, minimum height=1.5cm, align=center] (user) at (6,-4) {\textbf{End Users}\\Trust calibration\\Actionability};

% Arrows
\draw[->, very thick] (center) -- (ds);
\draw[->, very thick] (center) -- (bl);
\draw[->, very thick] (center) -- (reg);
\draw[->, very thick] (center) -- (user);

\end{tikzpicture}
\end{center}
\end{block}

\vspace{0.5cm}

%% AI LIFECYCLE INTEGRATION
\begin{block}{\Large Integrating Interpretability Across AI Lifecycle}

\begin{center}
\begin{tikzpicture}[scale=1.1, every node/.style={font=\Large}]
% Timeline
\draw[->, ultra thick] (0,0) -- (18,0);

% Phases
\node[draw, rectangle, rounded corners, fill=blue!30, minimum width=4cm, minimum height=2.5cm, align=center] at (2,0) {\textbf{Design Phase}\\[0.2cm]Choose models\\Define requirements\\Feature engineering};

\node[draw, rectangle, rounded corners, fill=green!30, minimum width=4cm, minimum height=2.5cm, align=center] at (8,0) {\textbf{Deployment}\\[0.2cm]Production architecture\\User interfaces\\Quality assurance};

\node[draw, rectangle, rounded corners, fill=orange!30, minimum width=4cm, minimum height=2.5cm, align=center] at (14,0) {\textbf{Monitoring}\\[0.2cm]Explanation quality\\Drift detection\\Compliance};

\end{tikzpicture}
\end{center}

\vspace{0.3cm}
\colorbox{blue!20}{\parbox{0.95\linewidth}{\large\textbf{Critical Insight:} Interpretability must be embedded from initial design, not retrofitted}}
\end{block}

\vspace{0.5cm}

%% CROSS-INDUSTRY LEARNINGS
\begin{block}{\Large Cross-Industry Learnings \& Applications}

\textbf{Key Observations from Case Studies:}

\begin{center}
\begin{tikzpicture}[scale=1.1, every node/.style={font=\large}]
% Domain boxes
\node[draw, rectangle, rounded corners, fill=red!20, minimum width=3.5cm, minimum height=2cm, align=center] at (0,6) {\textbf{Healthcare}\\GAMs, PDPs\\Clinical trust\\Sparse data challenges};

\node[draw, rectangle, rounded corners, fill=green!20, minimum width=3.5cm, minimum height=2cm, align=center] at (6,6) {\textbf{Finance}\\SHAP, EBMs\\Compliance\\Imbalanced data};

\node[draw, rectangle, rounded corners, fill=blue!20, minimum width=3.5cm, minimum height=2cm, align=center] at (12,6) {\textbf{LLMs/GenAI}\\Token attribution\\Opacity challenges\\Maturing field};

\node[draw, rectangle, rounded corners, fill=orange!20, minimum width=3.5cm, minimum height=2cm, align=center] at (3,2) {\textbf{Infrastructure}\\Root cause analysis\\ESG reporting\\Real-time constraints};

\node[draw, rectangle, rounded corners, fill=purple!20, minimum width=3.5cm, minimum height=2cm, align=center] at (9,2) {\textbf{HR/Education}\\Fairness\\Transparency\\Stakeholder communication};

\end{tikzpicture}
\end{center}

\vspace{0.5cm}

\textbf{Recurring Themes:}
\begin{itemize}[leftmargin=*]
\item \textbf{Early Integration Critical:} Projects defining interpretability goals early achieved better alignment
\item \textbf{Inherently Interpretable Models Preferred:} EBMs, GAMs, NAMs easier for stakeholders to understand
\item \textbf{Post-hoc Limitations:} SHAP and LIME struggled with stability in complex use cases
\item \textbf{Socio-Technical Challenge:} Success driven by process, communication, and alignment—not just technical execution
\end{itemize}
\end{block}

\vspace{0.5cm}

%% EVALUATION FRAMEWORK
\begin{block}{\Large Evaluation Framework}

\begin{center}
\begin{tikzpicture}[scale=1.1, every node/.style={font=\Large}]
% Main categories
\node[draw, rectangle, rounded corners, fill=blue!30, minimum width=5cm, minimum height=1.8cm, align=center] (quant) at (0,4) {\textbf{Quantitative}\\Fidelity, Stability\\Completeness};

\node[draw, rectangle, rounded corners, fill=green!30, minimum width=5cm, minimum height=1.8cm, align=center] (human) at (0,0) {\textbf{Human-Centered}\\Task performance\\Trust calibration\\Mental models};

\node[draw, rectangle, rounded corners, fill=orange!30, minimum width=5cm, minimum height=1.8cm, align=center] (co12) at (9,2) {\textbf{Co-12 Properties}\\Correctness\\Completeness\\Consistency\\Compactness\\Context\\Coherence};

% Arrows
\draw[->, very thick] (quant) -- (co12);
\draw[->, very thick] (human) -- (co12);

\end{tikzpicture}
\end{center}
\end{block}

\vspace{0.5cm}

%% RECOMMENDATIONS
\begin{block}{\Large Key Recommendations}
\begin{enumerate}[leftmargin=*, font=\Large\bfseries]
\item \textbf{Prioritize inherently interpretable models} over post-hoc explanations when performance allows
\item \textbf{Establish comprehensive evaluation frameworks} beyond traditional accuracy metrics
\item \textbf{Adapt interpretability strategies} to specific domain contexts and constraints
\item \textbf{Invest in stakeholder training} to properly evaluate and act on AI explanations
\item \textbf{Build interpretability into development workflows} from the design phase
\item \textbf{Develop communication infrastructure} (dashboards, visualizations, reports) for transparency
\end{enumerate}
\end{block}

\vspace{0.5cm}

%% CONCLUSION
\begin{block}{\Large Conclusion}
\begin{itemize}[leftmargin=*]
\item \textbf{Interpretability is not optional} for high-stakes AI deployment
\item \textbf{Success requires coordinated efforts} across technical, organizational, and human dimensions
\item \textbf{Organizations that embrace transparency as fundamental design principle} will achieve better compliance and unlock human-AI collaboration potential
\item \textbf{Proposed web portal initiative} to serve as centralized resource for interpretable AI implementation
\end{itemize}

\vspace{0.3cm}
\centering
\colorbox{green!30}{\parbox{0.9\linewidth}{\centering\Large\textbf{The future of AI depends on building systems humans can understand, trust, and effectively collaborate with}}}
\end{block}

\vspace{0.5cm}

\begin{block}{\Large References \& Contact}
\small
Full paper: \texttt{arXiv:2507.23535v1 [cs.LG]}\\
Vector Institute for Artificial Intelligence, Toronto\\
Contact: \texttt{https://vectorinstitute.ai}
\end{block}

\end{column}

\end{columns}
\end{frame}
\end{document}











% \documentclass[25pt, a0paper, portrait]{tikzposter}

% \usepackage{graphicx}
% \usepackage{booktabs}
% \usepackage{multirow}
% \usepackage{amsmath,amsfonts,amssymb}
% \usepackage{xcolor}
% \usepackage{enumitem}

% % Custom colors
% \definecolor{vectorblue}{RGB}{0, 102, 153}      % Vector Institute blue
% \definecolor{accentgreen}{RGB}{46, 204, 113}    % Green for highlights
% \definecolor{warningred}{RGB}{231, 76, 60}      % Red for warnings
% \definecolor{infoblue}{RGB}{52, 152, 219}       % Info blue

% % Theme settings
% \usetheme{Basic}
% \usecolorstyle{Default}

% \title{\textbf{Transparent AI: The Case for Interpretability and Explainability}}
% \author{Dhanesh Ramachandram, Himanshu Joshi, Judy Zhu, Dhari Gandhi, Lucas Hartman, Ananya Raval}
% \institute{Vector Institute for Artificial Intelligence, Toronto\\arXiv:2507.23535v1 [cs.LG]}

% \begin{document}

% \maketitle

% \begin{columns}
% \column{0.33}

% % Left column
% \block{Motivation \& Background}{
% \textbf{The Challenge:}
% \begin{itemize}[leftmargin=*]
%     \item AI systems increasingly inform \textbf{high-stakes decisions} in healthcare, finance, and public administration
%     \item \textcolor{warningred}{\textbf{Black box problem:}} Advanced models (deep neural networks) offer exceptional performance but lack transparency
%     \item Critical need to balance \textbf{model performance} and \textbf{interpretability}
% \end{itemize}

% \vspace{0.5cm}
% \textbf{The Example:}
% Consider a medical diagnostic tool with 2\% error rate but no explanation vs. a human physician with 15\% error rate but fully explainable decisions. \textbf{Trust requires understanding.}

% \vspace{0.5cm}
% \textbf{Key Sectors Requiring Transparency:}
% \begin{itemize}[leftmargin=*]
%     \item \textcolor{vectorblue}{\textbf{Healthcare:}} Clinicians need clarity on factors influencing AI predictions for cancer detection
%     \item \textcolor{vectorblue}{\textbf{Finance:}} Regulatory compliance demands transparent credit decisions (fair lending regulations)
%     \item \textcolor{vectorblue}{\textbf{Public Sector:}} Citizens expect clear explanations for decisions affecting their lives
% \end{itemize}
% }

% \block{Core Definitions}{
% \textbf{Interpretability vs. Explainability:}
% \begin{itemize}[leftmargin=*]
%     \item \textbf{Interpretability:} "How does the model function internally?"
%     \begin{itemize}
%         \item Addresses model's internal logic and decision-making processes
%         \item Inherent in transparent models (decision trees, linear regression)
%     \end{itemize}
%     \item \textbf{Explainability:} "Why did the model make a particular decision?"
%     \begin{itemize}
%         \item Provides clear reasons for specific decisions
%         \item Typically applied to complex black-box models (neural networks)
%     \end{itemize}
% \end{itemize}

% \vspace{0.5cm}
% \textbf{Taxonomy of Methods:}
% \begin{itemize}[leftmargin=*]
%     \item \textbf{Model-specific} vs. \textbf{Model-agnostic}
%     \item \textbf{Local} (individual predictions) vs. \textbf{Global} (overall behavior)
%     \item \textbf{Post hoc} (applied after training) vs. \textbf{Inherently interpretable} (built-in transparency)
% \end{itemize}
% }

% \block{Interpretability-Performance Trade-off}{
% \textbf{The Fundamental Challenge:}
% \begin{center}
% \Large
% \begin{tabular}{lcc}
% \toprule
% \textbf{Model Type} & \textbf{Performance} & \textbf{Interpretability} \\
% \midrule
% Deep Learning \& LLMs & \textcolor{accentgreen}{High} & \textcolor{warningred}{Low} \\
% Ensemble Models & \textcolor{accentgreen}{High} & \textcolor{warningred}{Medium} \\
% MLPs, SVMs & Medium & Medium \\
% Bayesian Models & Medium & \textcolor{accentgreen}{High} \\
% Linear Models & \textcolor{warningred}{Low} & \textcolor{accentgreen}{High} \\
% \bottomrule
% \end{tabular}
% \end{center}

% \vspace{0.3cm}
% \textbf{The Goal:} Bridge the gap between accuracy and transparency—develop high-performing models that are inherently interpretable.
% }

% \block{Stakeholder-Specific Guidelines}{
% \textbf{Four Key Stakeholder Groups:}
% \begin{enumerate}[leftmargin=*]
%     \item \textcolor{vectorblue}{\textbf{Data Scientists \& ML Specialists:}}
%     \begin{itemize}
%         \item Technical implementation of interpretability methods
%         \item Model selection and evaluation frameworks
%     \end{itemize}
%     \item \textcolor{vectorblue}{\textbf{Business Leaders \& Decision Makers:}}
%     \begin{itemize}
%         \item Strategic understanding of interpretability ROI
%         \item Governance frameworks and change management
%     \end{itemize}
%     \item \textcolor{vectorblue}{\textbf{Regulators \& Policy Makers:}}
%     \begin{itemize}
%         \item Compliance with transparency requirements
%         \item Regulatory frameworks and standards
%     \end{itemize}
%     \item \textcolor{vectorblue}{\textbf{End Users:}}
%     \begin{itemize}
%         \item Clear, understandable explanations
%         \item Trust and effective human-AI collaboration
%     \end{itemize}
% \end{enumerate}
% }

% \column{0.33}

% % Middle column
% \block{AI Development Workflow Integration}{
% \textbf{Design Phase:}
% \begin{itemize}[leftmargin=*]
%     \item \textcolor{infoblue}{\textbf{Problem Formulation:}} Identify interpretability requirements early
%     \item \textcolor{infoblue}{\textbf{Model Selection:}} Choose architecture balancing performance and transparency
%     \item \textcolor{infoblue}{\textbf{Data Strategy:}} Prepare data with interpretability in mind
% \end{itemize}

% \vspace{0.5cm}
% \textbf{Deployment Phase:}
% \begin{itemize}[leftmargin=*]
%     \item Generate explanations in production systems
%     \item Ensure explanations are timely and accessible
%     \item Validate explanation quality and usefulness
% \end{itemize}

% \vspace{0.5cm}
% \textbf{Monitoring \& Maintenance:}
% \begin{itemize}[leftmargin=*]
%     \item Continuous evaluation of interpretability quality
%     \item Feedback loops for explanation improvement
%     \item Regular audits and compliance checks
% \end{itemize}
% }

% \block{Standardized Reporting Framework}{
% \textbf{Essential Elements to Document:}
% \begin{itemize}[leftmargin=*]
%     \item \textbf{Model Architecture:} Structure, parameters, training details
%     \item \textbf{Interpretability Methods:} Techniques used, rationale for selection
%     \item \textbf{Explanation Quality:} Evaluation metrics, validation results
%     \item \textbf{Limitations:} Known constraints, edge cases, uncertainties
%     \item \textbf{Stakeholder Needs:} Addressed requirements, use case alignment
% \end{itemize}

% \vspace{0.5cm}
% \textbf{Case Study: Medical Imaging Diagnostic Support}
% \begin{itemize}[leftmargin=*]
%     \item Demonstrated successful implementation of reporting framework
%     \item Integration of interpretability into clinical workflows
%     \item Validation of explanation quality with domain experts
% \end{itemize}
% }

% \block{Evaluation Methods}{
% \textbf{Key Evaluation Dimensions:}
% \begin{enumerate}[leftmargin=*]
%     \item \textbf{Faithfulness:} Do explanations accurately reflect model behavior?
%     \item \textbf{Completeness:} Are explanations comprehensive enough?
%     \item \textbf{Stability:} Are explanations consistent across similar inputs?
%     \item \textbf{Understandability:} Can users comprehend the explanations?
%     \item \textbf{Usefulness:} Do explanations support user decision-making?
% \end{enumerate}

% \vspace{0.3cm}
% \textbf{Challenges:}
% \begin{itemize}[leftmargin=*]
%     \item Quantitative evaluation remains difficult
%     \item Human factors and cognitive limitations
%     \item Need for domain-specific evaluation criteria
% \end{itemize}
% }

% \block{Cross-Industry Applications}{
% \textbf{Key Findings Across Domains:}
% \begin{itemize}[leftmargin=*]
%     \item \textcolor{accentgreen}{\textbf{Healthcare:}} Interpretability enables clinicians to integrate AI with clinical judgment
%     \item \textcolor{accentgreen}{\textbf{Finance:}} Required for regulatory compliance (58\% of XAI studies lack adequate transparency)
%     \item \textcolor{accentgreen}{\textbf{Public Sector:}} Essential for public accountability and trust
% \end{itemize}

% \vspace{0.5cm}
% \textbf{Cross-Sector Observations:}
% \begin{itemize}[leftmargin=*]
%     \item Technical excellence alone is insufficient
%     \item Understanding stakeholder mental models is critical
%     \item Teams achieving meaningful outcomes invest in workflow integration
%     \item Context-specific explanations outperform generic approaches
% \end{itemize}
% }

% \column{0.33}

% % Right column
% \block{Challenges \& Limitations}{
% \textbf{Technical Challenges:}
% \begin{itemize}[leftmargin=*]
%     \item Complexity of modern deep learning architectures
%     \item Computational overhead of explanation generation
%     \item Scalability to large-scale systems
% \end{itemize}

% \vspace{0.5cm}
% \textbf{Interpretability-Performance Trade-off:}
% \begin{itemize}[leftmargin=*]
%     \item Often inverse relationship between accuracy and interpretability
%     \item Need for context-aware balance
%     \item Emerging research on inherently interpretable high-performing models
% \end{itemize}

% \vspace{0.5cm}
% \textbf{Human Factors:}
% \begin{itemize}[leftmargin=*]
%     \item Cognitive limitations in processing complex explanations
%     \item Varied stakeholder needs and backgrounds
%     \item Importance of explanation presentation format
% \end{itemize}

% \vspace{0.5cm}
% \textbf{Organizational Barriers:}
% \begin{itemize}[leftmargin=*]
%     \item Resource constraints and expertise gaps
%     \item Change management challenges
%     \item Lack of standardized frameworks
% \end{itemize}
% }

% \block{The Web Portal Initiative}{
% \textbf{Proposed Centralized Resource:}
% \begin{itemize}[leftmargin=*]
%     \item \textcolor{vectorblue}{\textbf{Interactive Technique Selection Tools:}}
%     \begin{itemize}
%         \item Guide practitioners to appropriate methods
%         \item Based on use case and requirements
%     \end{itemize}
%     \item \textcolor{vectorblue}{\textbf{Comprehensive Case Study Database:}}
%     \begin{itemize}
%         \item Real-world implementation examples
%         \item Cross-industry learnings and best practices
%     \end{itemize}
%     \item \textcolor{vectorblue}{\textbf{Regulatory Compliance Guidance:}}
%     \begin{itemize}
%         \item Sector-specific requirements
%         \item Compliance frameworks and checklists
%     \end{itemize}
% \end{itemize}
% }

% \block{Industry Adoption Roadmap}{
% \textbf{Key Components:}
% \begin{enumerate}[leftmargin=*]
%     \item \textbf{Training Programs:} Targeted education for technical and business stakeholders
%     \item \textbf{Pilot Projects:} Guided implementation in high-priority applications
%     \item \textbf{Knowledge Sharing:} Cross-functional learning and communities of practice
%     \item \textbf{Technical Workstreams:} Explore advanced challenges (LLMs, multi-modal systems)
%     \item \textbf{Leadership Initiatives:} Build strategic understanding among executives
%     \item \textbf{Continuous Improvement:} Regular feedback and iteration mechanisms
% \end{enumerate}

% \vspace{0.5cm}
% \textbf{Flexible Framework:}
% Adaptable to varying maturity levels, resource constraints, and sector-specific demands.
% }

% \block{Key Recommendations}{
% \textbf{For Organizations:}
% \begin{itemize}[leftmargin=*]
%     \item \textcolor{accentgreen}{✓} Assess current transparency practices
%     \item \textcolor{accentgreen}{✓} Identify high-priority applications requiring interpretability
%     \item \textcolor{accentgreen}{✓} Develop implementation roadmaps
%     \item \textcolor{accentgreen}{✓} Integrate interpretability as a core design principle
% \end{itemize}

% \vspace{0.5cm}
% \textbf{Critical Insights:}
% \begin{itemize}[leftmargin=*]
%     \item \textbf{Start early:} Interpretability should be built-in, not retrofitted
%     \item \textbf{Understand stakeholders:} Invest in understanding user mental models
%     \item \textbf{Context matters:} Tailor explanations to domain and use case
%     \item \textbf{Collaborative effort:} Requires coordination across organizational dimensions
% \end{itemize}
% }

% \block{Conclusions}{
% \textbf{The Imperative:}
% Interpretability has transformed from an academic consideration to a \textbf{fundamental requirement} for responsible AI deployment in critical decision-making roles.

% \vspace{0.5cm}
% \textbf{The Opportunity:}
% Organizations that embrace interpretability as a fundamental design principle will:
% \begin{itemize}[leftmargin=*]
%     \item Achieve better regulatory compliance
%     \item Unlock the full potential of human-AI collaboration
%     \item Build trust and accountability
% \end{itemize}

% \vspace{0.5cm}
% \textbf{The Future:}
% The future of AI depends on our collective ability to build systems that humans can \textbf{understand}, \textbf{trust}, and \textbf{effectively collaborate with}.

% \vspace{0.5cm}
% \textbf{The Journey:}
% The frameworks presented provide the foundation—\textcolor{vectorblue}{\textbf{the journey begins with commitment and action.}}
% }

% \block{Citation \& Resources}{
% \textbf{Paper:}
% Ramachandram, D., Joshi, H., Zhu, J., Gandhi, D., Hartman, L., \& Raval, A. (2025).\\
% \textit{Transparent AI: The Case for Interpretability and Explainability}.\\
% arXiv:2507.23535v1 [cs.LG]

% \vspace{0.5cm}
% \textbf{Key References:}
% \begin{itemize}[leftmargin=*]
%     \item Doshi-Velez \& Kim (2017): Evaluation and generalization in interpretable ML
%     \item Nauta et al. (2023): Systematic review on evaluating explainable AI
%     \item Rudin et al. (2022): Interpretable ML: Fundamental principles and grand challenges
%     \item Molnar (2020): Interpretable Machine Learning
% \end{itemize}
% }

% \end{columns}

% \end{document}


